% 论文排版样式（XeLaTeX + ctex）。字体配置由 build_paper_pdf.py 按本机环境注入：
% 优先采用优秀论文 A092 的宋体/黑体/Times New Roman 组合，缺少系统字体时回退 Fandol。
\usepackage[fontset=none]{ctex}
\usepackage{unicode-math}
@FONT_SETUP@
\usepackage{geometry}
\usepackage{indentfirst}
\usepackage{titlesec}
\usepackage{caption}
\usepackage{float}
\usepackage{setspace}
\usepackage{booktabs}
\usepackage{longtable}
\usepackage{graphicx}
\usepackage{fancyhdr}
\usepackage{amsmath,amssymb}
\usepackage{needspace}   % 表题与表体绑定，见 build_paper_pdf.py
\usepackage{fvextra}
\usepackage{seqsplit}

% 代码逐字环境必须覆盖中文注释、希腊字母与常用校验符号。
% 具体字体由构建脚本按本机可用字库注入，缺少宽字符字体时回退 DejaVu Sans Mono。
@CODE_FONT_SETUP@

% 页面：A4，四边留白不小于 2.5 cm
\geometry{a4paper, top=2.5cm, bottom=2.5cm, left=2.5cm, right=2.5cm}

% 图表前缀中文化
\renewcommand{\figurename}{图}
\renewcommand{\tablename}{表}

% 段落：小四号、单倍行距、首行缩进两字。
% 正文中文使用宋体，英文与数字使用 Times New Roman（字体配置见构建脚本）。
\setlength{\parindent}{2em}
\setstretch{1.0}
\setlength{\parskip}{0pt}

% 正文重点统一使用黑体粗体；若重点中含数学量，同时启用 \boldmath。
% 这比宋体粗体具有更清楚的层级，又不会放大字号。
\DeclareTextFontCommand{\textbf}{\bfseries\sffamily\boldmath}

% 页脚居中页码，无页眉横线（页眉不得出现任何身份信息）
\pagestyle{fancy}
\fancyhf{}
\fancyfoot[C]{\thepage}
\renewcommand{\headrulewidth}{0pt}
\fancypagestyle{plain}{\fancyhf{}\fancyfoot[C]{\thepage}\renewcommand{\headrulewidth}{0pt}}

% 标题层级。pandoc 以 --top-level-division=section 转换：
%   Markdown #   -> \section      论文主标题
%   Markdown ##  -> \subsection   一级标题（摘要 / 1 问题重述 / 2 数据审计 …）
%   Markdown ### -> \subsubsection 二级标题（1.1、1.2 …）
% 论文主标题为三号黑体；“一、问题重述”一类一级标题为小三黑体居中；
% “1.1 ……”二级标题为四号黑体左对齐。编号已经写在 Markdown 标题文本中，故不另生成编号。
\titleformat{\section}[block]{\centering\fontsize{16pt}{18pt}\selectfont\bfseries\sffamily}{}{0pt}{}
\titlespacing*{\section}{0pt}{8pt}{10pt}
\titleformat{\subsection}[block]{\centering\fontsize{15pt}{20pt}\selectfont\bfseries\sffamily}{}{0pt}{}
\titlespacing*{\subsection}{0pt}{10pt}{6pt}
\titleformat{\subsubsection}[block]{\fontsize{14pt}{18pt}\selectfont\bfseries\sffamily}{}{0pt}{}
\titlespacing*{\subsubsection}{0pt}{7pt}{3pt}
\titleformat{\paragraph}[block]{\fontsize{12pt}{16pt}\selectfont\bfseries\sffamily}{}{0pt}{}
\titlespacing*{\paragraph}{0pt}{5pt}{2pt}

% 图表题注：五号宋体/Times New Roman 加粗、居中、无缩进；图题在图下、表题在表上。
\captionsetup{justification=centering, font={small,bf}, labelfont=bf,
  singlelinecheck=true, skip=5pt}

% 摘要标题：三号黑体加粗，字间空一格，段前、段后各一行。
\newcommand{\paperabstractheading}{%
  \par\vspace{\baselineskip}%
  \begin{center}\fontsize{16pt}{16pt}\selectfont\bfseries\sffamily 摘\hspace{1em}要\end{center}%
  \vspace{\baselineskip}\par}

% 公式编号右对齐
\makeatletter
\def\tagform@#1{\maketag@@@{(#1\@@italiccorr)}}
\makeatother

% 附录源程序：允许在任意位置断行，避免长行溢出版心；字号压到 7.5pt 以控制篇幅
\fvset{breaklines=true, breakanywhere=true, fontsize=\fontsize{7.5pt}{9pt}\selectfont}
\DefineVerbatimEnvironment{Highlighting}{Verbatim}{breaklines=true, breakanywhere=true,
  commandchars=\\\{\}, fontsize=\fontsize{7.5pt}{9pt}\selectfont}

% 表格整体采用五号字和三线表风格。pandoc 生成的 longtable 本身使用 booktabs 横线，
% 这里关闭竖线并收紧列间距，使版面接近 A092 的表格密度。
\setlength{\tabcolsep}{5pt}
\renewcommand{\arraystretch}{1.05}
\AtBeginEnvironment{longtable}{\small\setlength{\LTpre}{3pt}\setlength{\LTpost}{3pt}}
\AtBeginEnvironment{tabular}{\small}
